% Zone „Das kennst du schon“ – aus bank/pyramide-kegel-kugel/zone.jsonl (zusammenbau v0.1)
\einheitenkopf[zone][Kennst du schon]{Das kennst du schon}
\setcounter{aufgabe}{0}
%% TODO Grundvorstellungs-Aufgabe als erste Hauptnummer der Zone (2.8) fehlt in der Bank

%% TODO Titel: Ich-kann-Satz fehlt in der Bank (Platzhalter: Kettenname) – Nr. 1
%% TODO Anweisung: ein Satz über der ersten Teilaufgabe fehlt in der Bank – Nr. 1
\begin{aufgabe}{Flächeninhalt von Quadrat, Rechteck, Dreieck (Seitenflächen) und Kreis (Grundkreis)}
%% TODO Darstellung für schwach fehlt in der Bank (pyramide-kegel-kugel-zone-f1-v1; Abschnitt „Für schwache Schüler“)
\swz{Quadrat mit der Seite $7$ cm – Flächeninhalt?}{}
%% TODO Darstellung für schwach fehlt in der Bank (pyramide-kegel-kugel-zone-f1-v3; Abschnitt „Für schwache Schüler“)
\swz{Kreis mit dem Radius $3$ cm – Flächeninhalt? Runde auf eine Stelle.}{}
\end{aufgabe}

%% TODO Titel: Ich-kann-Satz fehlt in der Bank (Platzhalter: Kettenname) – Nr. 2
%% TODO Anweisung: ein Satz über der ersten Teilaufgabe fehlt in der Bank – Nr. 2
\begin{aufgabe}{Volumen von Prisma und Zylinder als Grundfläche mal Höhe (der Bezugskörper für das Drittel) und Mantel des Zylinders}
%% TODO Darstellung für schwach fehlt in der Bank (pyramide-kegel-kugel-zone-f2-v1; Abschnitt „Für schwache Schüler“)
\swz{Prisma mit der Grundfläche $8$ cm² und der Höhe $6$ cm – Volumen?}{}
%% TODO Darstellung für schwach fehlt in der Bank (pyramide-kegel-kugel-zone-f2-v3; Abschnitt „Für schwache Schüler“)
\swz{Zylinder mit dem Radius $2$ cm und der Höhe $9$ cm – Volumen? Runde auf eine Stelle.}{}
\end{aufgabe}

%% TODO Titel: Ich-kann-Satz fehlt in der Bank (Platzhalter: Kettenname) – Nr. 3
%% TODO Anweisung: ein Satz über der ersten Teilaufgabe fehlt in der Bank – Nr. 3
\begin{aufgabe}{Satz des Pythagoras}
%% TODO Darstellung für schwach fehlt in der Bank (pyramide-kegel-kugel-zone-f3-v1; Abschnitt „Für schwache Schüler“)
\swz{Rechtwinkliges Dreieck, Katheten $3$ cm und $4$ cm – Hypotenuse?}{}
%% TODO Darstellung für schwach fehlt in der Bank (pyramide-kegel-kugel-zone-f3-v3; Abschnitt „Für schwache Schüler“)
\swz{Rechtwinkliges Dreieck, Katheten $1{,}2$ m und $0{,}9$ m – Hypotenuse?}{}
\end{aufgabe}

%% TODO Titel: Ich-kann-Satz fehlt in der Bank (Platzhalter: Kettenname) – Nr. 4
%% TODO Anweisung: ein Satz über der ersten Teilaufgabe fehlt in der Bank – Nr. 4
\begin{aufgabe}{Quadrieren und hoch drei, Wurzel und Kubikwurzel mit dem Taschenrechner, Rechnen mit π, Runden auf eine sinnvolle Stelle}
%% TODO Darstellung für schwach fehlt in der Bank (pyramide-kegel-kugel-zone-f4-v1; Abschnitt „Für schwache Schüler“)
\swz{$3^3$ – wie viel?}{}
%% TODO Darstellung für schwach fehlt in der Bank (pyramide-kegel-kugel-zone-f4-v3; Abschnitt „Für schwache Schüler“)
\swz{$1{,}5^3$ – wie viel?}{}
\end{aufgabe}

%% TODO Titel: Ich-kann-Satz fehlt in der Bank (Platzhalter: Kettenname) – Nr. 5
%% TODO Anweisung: ein Satz über der ersten Teilaufgabe fehlt in der Bank – Nr. 5
\begin{aufgabe}{Ein Drittel und vier Drittel einer Zahl}
%% TODO Darstellung für schwach fehlt in der Bank (pyramide-kegel-kugel-zone-f5-v1; Abschnitt „Für schwache Schüler“)
\swz{Ein Drittel von $27$?}{}
%% TODO Darstellung für schwach fehlt in der Bank (pyramide-kegel-kugel-zone-f5-v3; Abschnitt „Für schwache Schüler“)
\swz{Vier Drittel von $4{,}5$?}{}
\end{aufgabe}

%% TODO Titel: Ich-kann-Satz fehlt in der Bank (Platzhalter: Kettenname) – Nr. 6
%% TODO Anweisung: ein Satz über der ersten Teilaufgabe fehlt in der Bank – Nr. 6
\begin{aufgabe}{Formel aus der Formelsammlung entnehmen und nach einer Größe umstellen (Volumenformel mit Bruchfaktor nach der Höhe)}
%% TODO Darstellung für schwach fehlt in der Bank (pyramide-kegel-kugel-zone-f6-v1; Abschnitt „Für schwache Schüler“)
\swz{Rechteck: $A = a \cdot b$ mit $A = 42$ m² und $a = 6$ m. Wie lang ist $b$?}{}
%% TODO Darstellung für schwach fehlt in der Bank (pyramide-kegel-kugel-zone-f6-v3; Abschnitt „Für schwache Schüler“)
\swz{Prisma: $V = G \cdot h$ mit dem Volumen $V = 37{,}5$ dm³ und der Grundfläche $G = 7{,}5$ dm². Wie hoch ist das Prisma?}{}
\end{aufgabe}

%% TODO Titel: Ich-kann-Satz fehlt in der Bank (Platzhalter: Kettenname) – Nr. 7
%% TODO Anweisung: ein Satz über der ersten Teilaufgabe fehlt in der Bank – Nr. 7
\begin{aufgabe}{Radius aus Durchmesser}
%% TODO Darstellung für schwach fehlt in der Bank (pyramide-kegel-kugel-zone-f7-v1; Abschnitt „Für schwache Schüler“)
\swz{Durchmesser $18$ cm – Radius?}{}
%% TODO Darstellung für schwach fehlt in der Bank (pyramide-kegel-kugel-zone-f7-v3; Abschnitt „Für schwache Schüler“)
\swz{Durchmesser $45$ mm – Radius in cm?}{}
\end{aufgabe}

%% TODO Titel: Ich-kann-Satz fehlt in der Bank (Platzhalter: Kettenname) – Nr. 8
%% TODO Anweisung: ein Satz über der ersten Teilaufgabe fehlt in der Bank – Nr. 8
\begin{aufgabe}{Volumeneinheiten cm³ → dm³ = l mit 1000, m³; Masse aus Volumen und Dichte}
%% TODO Darstellung für schwach fehlt in der Bank (pyramide-kegel-kugel-zone-f8-v1; Abschnitt „Für schwache Schüler“)
\swz{$3\,000$ cm³ – wie viel Liter?}{}
%% TODO Darstellung für schwach fehlt in der Bank (pyramide-kegel-kugel-zone-f8-v3; Abschnitt „Für schwache Schüler“)
\swz{$750$ cm³ – wie viel Liter?}{}
\end{aufgabe}

%% TODO Titel: Ich-kann-Satz fehlt in der Bank (Platzhalter: Kettenname) – Nr. 9
%% TODO Anweisung: ein Satz über der ersten Teilaufgabe fehlt in der Bank – Nr. 9
\begin{aufgabe}{Termumformung mit Variablen}
%% TODO Darstellung für schwach fehlt in der Bank (pyramide-kegel-kugel-zone-f9-v1; Abschnitt „Für schwache Schüler“)
\swz{$(2a)^2$ – vereinfache.}{}
%% TODO Darstellung für schwach fehlt in der Bank (pyramide-kegel-kugel-zone-f9-v3; Abschnitt „Für schwache Schüler“)
\swz{$(0{,}5b)^2$ – vereinfache.}{}
\end{aufgabe}

%% TODO Titel: Ich-kann-Satz fehlt in der Bank (Platzhalter: Kettenname) – Nr. 10
%% TODO Anweisung: ein Satz über der ersten Teilaufgabe fehlt in der Bank – Nr. 10
\begin{aufgabe}{Radius aus Durchmesser – Fehler finden}
\swz[3]{Kreis mit dem Durchmesser $16$ cm. Lena rechnet: \rechnung{A &= \pi \cdot 16^2 \\ A &\approx 804{,}2 \text{ cm}^2} Finde den Fehler.}{}
\end{aufgabe}

%% TODO Titel: Ich-kann-Satz fehlt in der Bank (Platzhalter: Kettenname) – Nr. 11
%% TODO Anweisung: ein Satz über der ersten Teilaufgabe fehlt in der Bank – Nr. 11
\begin{aufgabe}{Radius aus Durchmesser – selbst rechnen}
%% TODO Darstellung für schwach fehlt in der Bank (pyramide-kegel-kugel-zone-f7-v6; Abschnitt „Für schwache Schüler“)
\swz{Kreis mit dem Durchmesser $14$ cm – Flächeninhalt? Runde auf eine Stelle.}{}
\end{aufgabe}
